\section{Evaluation Methodology}

\subsection{Representation and engine tests}

Evaluation begins below gameplay. Layer-off tests require bit-identical base
logits, legal choices, and Card2Vec outputs. The collision census classifies
every legacy duplicate by public semantic delta, selected-action involvement,
and available successor witness. Targeted engine cases cover full-bench attack
requirements, deck-zero or opponent-hand-zero conditions, automatic effect
ordering, simultaneous-knockout prompt order, ordinary/ex/Mega-ex Prize counts,
and visible Prize modifiers. A legality predictor may flag a discrepancy but
cannot alter the engine list.

Information-set metamorphic tests hold the acting player's sanitized state,
history, and legal options fixed while changing privileged opponent hands, deck
order, or unrevealed Prizes. Observation hashes, option features, base logits,
candidate logits, and any public-belief proposal must remain identical. The
realized future may change after action execution; the pre-transition decision
may not depend on that hidden truth.

\subsection{Paired policy comparisons}

The smallest reported gameplay comparison is BO250: 125 seed-matched pairs,
with candidate and reference swapping seats inside each pair. BO250 is a
diagnostic, not a conclusive promotion gate. Reports include wins, draws,
losses, paired intervals, action disagreement, severe rule-error rate, residual
coverage, and which bounded contribution changed each action.

Formal evaluation uses a fixed opponent roster and 50 games per opponent deck,
split 25/25 by candidate seat. With the current 18-deck roster this yields 900
games. We report skill-weighted rate, its receipt-defined lower bound, minimum
opponent rate, tier strata, and every pipeline criterion. A statistically
favorable aggregate does not erase a failed rating simulation, historical
non-regression check, or identity audit. Wilson-style binomial intervals are
useful descriptive summaries~\cite{wilson1927probable}, but the receipt-defined
gate remains the operational criterion.

The Prize-plan study has a separate preregistered structure. Eligible legacy,
H1, H3, and later H6 arms each require 8,192 seed-matched, seat-swapped games
when capacity permits. Required outputs include W/D/L, Prize conversion,
matchup and public-Prize strata, late closeout, deck-out, action entropy,
calibration, coverage, AWR effective sample size, clipping, and action
disagreement. No such completed multi-arm gameplay receipt exists at the paper
snapshot.

\subsection{Evidence tiers}

We distinguish five tiers. \Measured\ is a completed numerical artifact with
its experimental unit. \Verified\ is an integration, identity, or training
receipt that does not itself prove game strength. \Configured\ is a future
protocol or active design. \Historical\ is context from a different checkpoint,
deck, or uncontrolled public environment. \NotMeasured\ means no result is
available, including a waived test. These labels remain adjacent to each table
row so a detached number cannot silently acquire a stronger interpretation.
